\begin{tikzpicture}[x=5.2cm,y=1.05cm,
  axis/.style={-{Latex[length=1.4mm]},semithick},
  guide/.style={thin,black!45,densely dotted},
  lab/.style={font=\scriptsize,align=center}]
\draw[axis] (0,0) -- (0,3.22) node[above,font=\scriptsize] {flux [arb.]};
\draw[axis] (-0.015,0) -- (1.12,0) node[below,font=\scriptsize] {$\xi$};
\foreach \xx in {0.25,0.5,0.6667,0.75,1.0}{\draw (\xx,0.045) -- (\xx,-0.045) node[below,font=\tiny] {\xx};}
\draw[black!40] (0,1.0) -- (1,0);
\draw[black!40,densely dashed] (0,0) -- (1,1.0);
\draw[black!62] (0,2.0) -- (1,0);
\draw[black!62,densely dashed] (0,0) -- (1,1.0);
\draw[thick] (0,3.0) -- (1,0);
\draw[thick,densely dashed] (0,0) -- (1,1.0);
\foreach \x/\y/\txt/\pos in {0.5000/0.5000/{$A=0$}/left,0.6667/0.6667/{$r^+=2r^-$}/right,0.7500/0.7500/{$r^+=3r^-$}/right}{
  \draw[guide] (\x,0) -- (\x,\y);
  \fill (\x,\y) circle (0.8pt);
  \node[font=\scriptsize,\pos] at (\x,\y+0.10) {\txt};
}
\node[lab,align=left] at (0.18,2.72) {forward\\$r^+(1-\xi)$};
\node[lab,align=left] at (0.84,1.15) {reverse\\$r^-\xi$};
\draw[-{Latex[length=1.2mm]},semithick] (0.49,1.72) -- (0.67,0.77)
  node[midway,right,font=\scriptsize] {$\mathcal A\uparrow$};
\node[lab] at (0.62,2.35) {$\xi_\eq=r^+/(r^++r^-)$};
\end{tikzpicture}
\caption{정지점 유도 개선안. 정방향 플럭스 $r^+(1-\xi)$ 와 역방향 플럭스 $r^-\xi$ 의 교점이
$\xi_\eq$(식~\eqref{eq:logisticsolve})이다. 좌표는 $r^-=1$ 로 두고 $r^+/r^-=\exp(\mathcal A/RT)$ 를
T5\_calc.py 로 수치 평가했다. $\mathcal A=0$ 이면 $\xi_\eq=\tfrac12$, $\mathcal A/RT=\ln2$ 이면
$\xi_\eq=2/3$, $\mathcal A/RT=\ln3$ 이면 $\xi_\eq=3/4$ 이며 detailed balance(식~\eqref{eq:db})가 교점 이동을 정한다.}
\label{fig:flux}
